\documentclass[10pt,a4paper]{amsart}
\usepackage[margin=19mm]{geometry}
\usepackage{amsmath,amssymb,mathtools}
\usepackage[scr=rsfs,cal=euler]{mathalfa}
\usepackage{xfrac}
\usepackage[unicode,hidelinks,bookmarks=false]{hyperref}
\newcommand{\ie}{i.e.\ }
\newcommand{\cf}{cf.\ }
\newcommand{\resp}{respectively\ }
\newcommand{\Subsection}{\subsection}
\providecommand{\nobreakdash}{\nobreak\hbox{-}}
\setlength{\emergencystretch}{2em}
\multlinegap0pt
\usepackage{tcolorbox}
\usepackage{tikz}
\usepackage{bm}
\newtheorem{lemma}{Lemma}
\newtheorem{proposition}{Proposition}
\newtheorem{corollary}{Corollary}
\usepackage{cleveref}
\crefname{lemma}{Lemma}{Lemmas}
\crefname{proposition}{Proposition}{Propositions}
\crefname{corollary}{Corollary}{Corollarys}
\newcommand{\botmap}{\operatorname{bot}}
\newcommand{\cT}{\mathcal{T}}
\newcommand{\topmap}{\operatorname{top}}
\title{Counting Boolean antichains}
\author{Julien Garber, Lina S. Goltermann, Daniel Horiatakis, Dennis K\"onig, Tal Gottesman}
\date{Source: arXiv:2608.27126v2 (CC BY 4.0)}
\begin{document}
\maketitle

\noindent\emph{Source and licence.} Counting Boolean antichains. Version arXiv:2608.27126v2 (CC BY 4.0), distributed by arXiv under the Creative Commons Attribution 4.0 International licence, http://creativecommons.org/licenses/by/4.0/, as recorded in the arXiv licence metadata for this exact version and shown on the abstract page https://arxiv.org/abs/2608.27126v2. Full licence notice: \texttt{licenses/arxiv-2608.27126-CC-BY-4.0.txt}.\par\smallskip

\noindent\textit{Editorial scope.} Counted unit: the article's proposition ValidAreBoolean (source0.tex lines 1182--1184). The article's complete original proof of it is delivered with it (source0.tex lines 1185--1234, one \texttt{proof} environment of 50 lines). No mathematics is written, repaired or re-derived: the statement and the proof are the article's own lines, in the article's own notation, and no retained line is substituted or re-flowed.  Retained uncounted as the source-local context the proof needs: lines 150--152; lines 663--665; lines 779--785; lines 878--880; lines 891--897; lines 936--938; lines 954--956; lines 960--962; lines 1096--1098; lines 1138--1140; lines 1143--1145; lines 1150--1153; lines 1170--1175.  Nothing was omitted from any retained span, so the package carries no omission record.  No retained line contains a citation, so the wrapper carries no bibliography and no bibliography entry.  No retained line contains a figure environment, an image-inclusion command or drawing code, so no image is delivered and the package declares no asset dependency.  Editorial formatting only, in addition: an AMS \texttt{amsart} preamble that restates the article's own declarations and macros used by the retained lines, namely the environments \texttt{lemma}, \texttt{proposition}, \texttt{corollary} on separate counters, together with the standard packages those lines need.  The retained environments print in this document as Lemma 1, Proposition 1, Corollary 1 in order of appearance, whereas the article prints the same items with the numbering of its own full document.  The complete native source of the article as distributed in the arXiv e-print for this exact version is bundled unchanged as \texttt{sources/208689/source0.tex}.
\begin{equation}\label{eq: intersectivity}
    (\wedge S) \vee (\wedge S') = \wedge (S\cap S').
\end{equation}

\begin{equation}\label{eq:coverRel}
(\dots((xy)z)\dots) \lessdot (\dots (x(yz))\dots).
\end{equation}

\begin{lemma}[Banana Lemma\footnote{This name stems from the habit of the first four authors of referring to the equivalence classes under \emph{q} as "bananas". We found the name catchy and decided to keep it.}]\label{lem:banana}
    Let $a$ and $b$ be elements in $\cT_{n}$ such that $a \lessdot b$. Let $z$ be the subexpression on the right of the reparenthesization that transforms $a$ into $b$ as in \Cref{eq:coverRel}.
    \begin{enumerate}
        \item If $x_{n+1} \notin z$, then the element of height $i$ in $q^{-1}(\{a\})$ is covered by the element of height $i$ in $q^{-1}(\{b\})$ for every possible height $i$. In particular, both classes have the same cardinality.
        \item If $x_{n+1} \in z$, let $j$ be the number of closing parentheses after $z$. Then the element of height $i$ in $q^{-1}(\{a\})$ is covered by the element of height $i$ in $q^{-1}(\{b\})$ if $i < j$ and by the element of height $i+1$ otherwise. In particular, $|q^{-1}(\{b\})| = |q^{-1}(\{a\})| + 1$.
    \end{enumerate}
\end{lemma}

\begin{proposition}\label{prop: LatticeCongruence}
    The map $q$ defines a lattice congruence on $\cT_{n+1}$.
\end{proposition}

\begin{corollary}\label{cor: bot and top meet and join}
    Let $a,b \in \cT_{n+1}$. Then
    \begin{enumerate}
        \item $\topmap(a)\wedge \topmap(b) = \topmap(a\wedge b)$, $\topmap(a)\vee \topmap(b) = \topmap(a\vee b)$,
        \item $\botmap(a)\wedge \botmap(b) = \botmap(a\wedge b)$ and $\botmap(a)\vee \botmap(b) = \botmap(a\vee b)$.
    \end{enumerate}
\end{corollary}

\begin{lemma}\label{lem: closing parentheses in equivalence class}
    Let $v\lessdot w \in \cT_{n+1}$ and let $z$ be the subexpression on the right of the reparenthesization that transforms $v$ into $w$ as in \Cref{eq:coverRel}. If $x_{n+2}\notin z$, then $I_v = I_w$. If $x_{n+2}\in z$, then $I_w\subsetneq I_v$. Additionally, if $v$ and $w$ are elements of the same equivalence class, we have $I_v = I_w \cup \{k_w\}$.
\end{lemma}

\begin{corollary}\label{cor: closing parenthesis moved to the right}
    Let $v \leq w \in \cT_{n+1}$, then $I_w \subseteq I_v$.
\end{corollary}

\begin{lemma}\label{lem: meet height is union}
    Let $v,w \in \cT_{n+1}$. Then we have $I_{v \wedge w} = I_v \cup I_w$.
\end{lemma}

\begin{proposition}\label{prop: BooleanAreValid}
    Let $C$ be a Boolean antichain of size $k\leq n$ in $\cT_{n+1}$. Then $C$ is an almost Boolean antichain of the form $C=\{f_n(c_1),\dots,f_n(c_{k-1}),w\}$, and there exists at most one $i \in [1, k-1]$ with $k_w \in I_{f_n(c_i)}$.
\end{proposition}

\begin{equation}\label{eq: pi of C Boolean}
    \pi_i^\ell(C) \text{ is a Boolean antichain for all } i\leq \ell.
\end{equation}

\begin{equation}\label{eq: I_q_i}
    \text{for } d\in \cT_\ell \text{ we have }I_{f_\ell(d)} = I_d \text{ and } I_{q_\ell(d)}\subseteq I_d
\end{equation}

\begin{lemma}\label{lem: backtracking}
    Let $C$ be an almost Boolean antichain in $T_\ell$. For a fixed $c \in C$ and a fixed $j \in I_c$ there exists an integer $m \leq \ell$ such that
    $\pi_m^\ell(c) = w(\pi_m^\ell(C))$ and $j \in [k_{\pi_m^\ell(c)} + 1, m - 1] \subseteq I_{\pi_m^\ell(c)}$.
\end{lemma}

\begin{lemma}\label{lem: forwarding k}
    Let $C$ be an almost Boolean antichain in $T_\ell$, $c \in C$ and $j \in I_c$. Then we have
    \[
        k_{\pi_{n(c,j)}^\ell(c)} \notin I_c.
    \]
\end{lemma}

\begin{proposition}\label{prop: ValidAreBoolean}
    Let $C=\{f_n(c_1),\dots,f_n(c_{k-1}),w\}$ be an almost Boolean antichain of size $k\leq n$ in $\cT_{n+1}$. Suppose that there exists at most one $i \in [1,k-1]$ with $k_w \in I_{f_n(c_i)}$. Then $C$ is a Boolean antichain.
\end{proposition}

\begin{proof}
    Let $S$ and $S'$ be subsets of $ C$. As $C$ is an almost Boolean antichain, \eqref{eq: intersectivity} is fulfilled for $\{c_1, \dots, c_{k-1}, q_n(w)\}$. Since $q$ is a lattice congruence (see \Cref{prop: LatticeCongruence}), the equality $[(\wedge S) \vee (\wedge S')] = [\wedge (S\cap S')]$ holds. By definition, we have $f_n(c_i) = \text{top}(f_n(c_i))$. If $w$ is in at most one of $S, S'$, one can use \Cref{lem:banana} and \Cref{cor: bot and top meet and join} to show that $(\wedge S) \vee (\wedge S')$ and $\wedge (S\cap S')$ have the same height. To do so in general, we compute $k_{\wedge (S\cap S')}$ and argue that $k_{\wedge (S\cap S')}\not\in I_{(\wedge S) \vee (\wedge S')}$, in order to apply \Cref{lem: closing parentheses in equivalence class}. The proof is divided into four steps.

    \noindent{\sf Step 1.}
    We first construct a sequence $(a_i)_{1\leq i\leq r}$ of elements in $C$ and a sequence $(d_i)_{0\leq i\leq r}$ of corresponding elements in different $\pi_{j_i}(C)$. To do this, we iteratively use the process described in \Cref{lem: backtracking} to project the latest $d_i$ onto a smaller Tamari lattice $\cT_m$, where it is a non-top element of its equivalence class. 
    If $i\geq 0$ we denote the index $k_{\pi_m(a_i)}$ by $\overline{k_{d_i}}$ where $m\leq n +1$ will be clear from the context. For $i = -1$, we set $\overline{k_{d_{-1}}} = k_{w}+1$.
    
    First, let $d_0$ be $w$. By the assumption on $C$, there exists at most one element $d_1 \in C$ with $k_w \in I_{d_1}$ and we define $a_1:=d_1$.
    Now suppose we have $a_1, \dots, a_{i}$ and $d_0, \dots, d_{i}$ satisfying $\overline{k_{d_{i-1}}}\in I_{d_{i}}$, and $\pi_{n(d_{i-1}, \overline{k_{d_{i-2}}})}(a_{i}) = d_{i}$.
    %By \ref{lem: backtracking}, there exists $n(d_1, k_w)\leq n+1$. We know by \eqref{eq: pi of C Boolean} that $\pi_{n(d_1, k_w)}(C)$ is again a Boolean antichain. Thus, we know by \Cref{prop: BooleanAreValid} that there exists at most one element $d_2 \in \pi_{n(d_1, k_w)}(C)$ with $k_{\pi_{n(d_1, k_w)}(d_1)} \in I_{d_2}$. By construction we know that there exists $a_2 \in C$ with $\pi_{n(d_1, k_w)}(a_2) = d_2$. This $a_2$ is well-defined: If there exists another $a \in C$ with $\pi_{n(d_1, k_w)}(a) = d_2$, then in particular we have $[\pi_{n(d_1, k_w)+1}(a_2)] = [\pi_{n(d_1, k_w)+1}(a)]$. However, we then have two comparable elements in an antichain, yielding a contradiction.
    We construct $d_{i+1}$ and $a_{i+1}$. By \Cref{lem: backtracking}, there exists $n(d_i,\overline{k_{d_{i-1}}}) < n(d_{i-1},\overline{k_{d_{i-2}}})$ such that \[\pi_{n(d_i, \overline{k_{d_{i-1}}})}^{n(d_{i-1}, \overline{k_{d_{i-2}}})}(d_i) = w(\pi_{n(d_{i}, \overline{k_{d_{i-1}}})}(C)).\]
    By \eqref{eq: pi of C Boolean}, $\pi_{n(d_i,\overline{k_{d_{i-1}}})}(C)$ is a Boolean antichain. Hence, by \Cref{prop: BooleanAreValid}, there exists at most one $d_{i+1} \in \pi_{n(d_i,\overline{k_{d_{i-1}}})}(C)$ with $\overline{k_{d_i}} \in I_{d_{i+1}}$. If it exists, then there exists $a_{i+1} \in C$ with $\pi_{n(d_i, \overline{k_{d_{i-1}}})}(a_{i+1}) = d_{i+1}$. This $a_{i+1}$ is unique: If there exists another $a \in C$ such that $\pi_{n(d_i, \overline{k_{d_{i-1}}})}(a) = d_{i+1}$, then the two elements $\pi_{n(d_i, \overline{k_{d_{i-1}}})+1}(a_{i+1})$ and $\pi_{n(d_i, \overline{k_{d_{i-1}}})+1}(a)$ are in the same equivalence class. However, this gives two comparable elements in an antichain, so they must be equal. Applying this argument repeatedly yields $a = a_{i+1}$.
    
    We stop this process as soon as either there exists no $d_{r+1}$ such that $\overline{k_{d_r}} \in I_{d_{r+1}}$ holds, or $a_{r+1} \notin S \cap S'$. Since $n(d_i,\overline{k_{d_{i-1}}})$ strictly decreases in every iteration, one of these two cases occurs after a finite number of iterations.  
    Moreover, we have by \Cref{lem: backtracking} that
    \begin{equation}\label{eq: k_d_j}
        \overline{k_{d_j}} \in [\overline{k_{d_{j+1}}}+1,n(d_{j+1},\overline{k_{d_j}})-1] \text{ for all } 0 \leq j \leq r-1.
    \end{equation}

    \noindent{\sf Step 2.} We now show that for each $i$ there exists at most one $a \in C$ with $\overline{k_{d_i}} \in I_a$ and that, if it exists, then $i\leq r$, $d_{i+1}$ exists and $a$ must be $a_{i+1}$. In particular, this implies $\overline{k_{d_i}} \in I_{a_{i+1}}$.

    We proceed by induction on $i$. By the assumption of the theorem, the initial step holds, in which $a_1=d_1$ is in $C$. For the induction step, let $a \in C$ such that $\overline{k_{d_i}} \in I_a$. By \Cref{lem: backtracking}, there exists $n(a,\overline{k_{d_i}})$ such that
    \begin{equation}\label{eq: I_a_overline}
        \overline{k_{d_i}} \in [\overline{k_a}+1,n(a,\overline{k_{d_i}})-1] \subseteq I_{\pi_{n(a,\overline{k_{d_i}})}(a)}.    
    \end{equation}
    Suppose that $n(a,\overline{k_{d_i}}) > n(d_i,\overline{k_{d_{i-1}}})$, then \eqref{eq: k_d_j} and \eqref{eq: I_a_overline} yield that
    $$
        \overline{k_a}+1 \leq \overline{k_{d_i}} < \overline{k_{d_{i-1}}} < n(d_i,\overline{k_{d_{i-1}}}) \leq n(a,\overline{k_{d_i}})-1.
    $$
    We use \eqref{eq: I_a_overline} and \eqref{eq: I_q_i} to see that $\overline{k_{d_{i-1}}} \in I_a$. Since $\overline{k_{d_i}} \notin I_{a_i}$ holds by \Cref{lem: forwarding k}, but $\overline{k_{d_i}} \in I_a$ by assumption, $a$ and $a_i$ must be distinct. But by the induction hypothesis, $a_i$ is the only element of $C$ with $\overline{k_{d_{i-1}}}$ in its closing index set, yielding a contradiction. This shows that $n(a,\overline{k_{d_i}}) \leq n(d_i,\overline{k_{d_{i-1}}})$. Let $a'$ in $\cT_{n(d_i,\overline{k_{d_{i-1}}})}$ be such that $a' = \pi_{n(d_i,\overline{k_{d_{i-1}}})}(a)$. Then \eqref{eq: I_q_i} and \eqref{eq: I_a_overline} show that the integer $\overline{k_{d_i}}$ lies in $I_{a'}$. Thus, there exists an element in $\cT_{n(d_i,\overline{k_{d_{i-1}}})}$ with $\overline{k_{d_i}}$ in its index set. But since there exists at most one such element, namely $d_{i+1}$, we have $a'=d_{i+1}$. By the same argument as in {\sf Step 1}, $a$ is equal to $a_{i+1}$.

    \noindent{\sf Step 3.} We now compute $k_{\wedge S\cap S'}$.
    Let $(a_i)_{1\leq i\leq r}$ and $(d_i)_{0 \leq i\leq r}$ be the sequences constructed in {\sf Step 1}. By {\sf Step 2}, if there exists no $d_{r+1}$, then there also exists no $a \in C$ with $\overline{k_{d_r}} \in I_a$. On the other hand, if $d_{r+1}$ exists and $a_{r+1}\notin S\cap S'$, then $a_{r+1}$ is the only element of $C$ such that $\overline{k_{d_r}} \in I_{a_{r+1}}$. In both cases, \Cref{lem: meet height is union} can be applied to obtain $\overline{k_{d_r}} \notin I_{\wedge(S\cap S')}$. Recall that
    $$
    a_r,a_{r-1},\dots,a_1,w \in S\cap S'.
    $$
    Hence, \Cref{lem: meet height is union}, \Cref{lem: backtracking} and \eqref{eq: I_q_i} yield that
    \begin{align*}
        &[\overline{k_{d_r}}+1,n(d_r,\overline{k_{d_{r-1}}})-1] \cup [\overline{k_{d_{r-1}}}+1,n(d_{r-1},\overline{k_{d_{r-2}}})-1]\; \cup \\ &\dots \cup [\overline{k_{d_1}}+1,n(d_1,k_w)-1] \cup 
    [k_w+1,n]\subseteq I_{\wedge(S\cap S')}.
    \end{align*}
    Together with \eqref{eq: k_d_j}, this gives $j \in I_{\wedge(S\cap S')}$ for all $\overline{k_{d_r}}<j\leq n$. Hence, it follows that 
    $$
        k_{\wedge (S \cap S')} = \overline{k_{d_r}}.
    $$
    
    \noindent{\sf Step 4.} As in {\sf Step 3}, we observe that either no element of $C$ contains $\overline{k_{d_{r}}}$ in its closing index set or $a_{r+1}$ exists and $a_{r+1}\notin S\cap S'$. In the latter case, we can assume that $a_{r+1} \notin S$. Hence, no element in $S$ has $\overline{k_{d_r}}=k_{\wedge (S \cap S')}$ in its closing index set. By \Cref{lem: meet height is union}, the closing index set $I_{\wedge S}$ also does not contain $k_{\wedge (S \cap S')}$. Moreover, by \Cref{cor: closing parenthesis moved to the right}, we have 
    \[I_{(\wedge S)\vee (\wedge S')} \subseteq I_{\wedge S},\]
    so $k_{\wedge (S \cap S')}$ is not in $I_{(\wedge S)\vee (\wedge S')}$. Recall that $(\wedge S)\vee (\wedge S') \leq \wedge(S\cap S')$ holds and that $(\wedge S)\vee (\wedge S')$ and $ \wedge(S\cap S')$ lie in the same equivalence class. By \Cref{lem: closing parentheses in equivalence class}, any element lying strictly below $\wedge (S\cap S')$ contains $k_{\wedge(S\cap S')}$ in its closing index set. Hence $\wedge (S\cap S')$ and $(\wedge S) \vee (\wedge S')$ are equal, showing that $C$ satisfies \eqref{eq: intersectivity}. Thus, $C$ is a Boolean antichain. 
\end{proof}

\end{document}
